\documentclass[11pt,a4paper]{article}

\usepackage[utf8]{inputenc}
\usepackage[T1]{fontenc}
\usepackage[french]{babel}
\usepackage[a4paper,margin=2.2cm]{geometry}
\usepackage{amsmath,amssymb,amsthm,mathtools}
\usepackage{lmodern}
\usepackage{enumitem}
\usepackage{hyperref}
\usepackage{xcolor}
\usepackage{fancyhdr}
\usepackage{stmaryrd}
\usepackage{mathrsfs}
\usepackage{array}

\hypersetup{
    colorlinks=true,
    linkcolor=black,
    urlcolor=blue,
    citecolor=black
}

\setlength{\parindent}{0pt}
\setlength{\parskip}{0.6em}

\pagestyle{fancy}
\fancyhf{}
\fancyhead[L]{Algèbre linéaire --- Chapitre 10}
\fancyhead[R]{Déterminants}
\fancyfoot[C]{\thepage}

% Environnements
\newtheorem{definition}{Définition}[section]
\newtheorem{theoreme}[definition]{Théorème}
\newtheorem{proposition}[definition]{Proposition}
\newtheorem{corollaire}[definition]{Corollaire}
\newtheorem{lemme}[definition]{Lemme}
\theoremstyle{remark}
\newtheorem{remarque}[definition]{Remarque}
\newtheorem{exemple}[definition]{Exemple}
\newtheorem*{methode}{Méthode}
\newtheorem*{idee}{Idée}

% Commandes
\newcommand{\R}{\mathbb{R}}
\newcommand{\C}{\mathbb{C}}
\newcommand{\K}{\mathbb{K}}
\newcommand{\N}{\mathbb{N}}
\newcommand{\Z}{\mathbb{Z}}
\newcommand{\Vect}{\mathrm{Vect}}
\newcommand{\Ker}{\mathrm{Ker}}
\newcommand{\Img}{\mathrm{Im}}
\newcommand{\rg}{\mathrm{rg}}
\newcommand{\id}{\mathrm{Id}}
\newcommand{\tr}{\mathrm{tr}}
\newcommand{\Mnn}{M_n(\K)}
\newcommand{\GLn}{GL_n(\K)}

\begin{document}

\begin{center}
    {\LARGE \textbf{Chapitre 10 --- Déterminants}}\\[0.4em]
    {\large Version complète et détaillée}
\end{center}

\hrule
\vspace{1em}

\tableofcontents

\newpage

\section{Pourquoi introduire le déterminant ?}

Le déterminant est une fonction associée aux matrices carrées qui joue un rôle central en algèbre linéaire. Il permet notamment :
\begin{itemize}
    \item de caractériser l’inversibilité d’une matrice ;
    \item de détecter la dépendance linéaire de colonnes ou de lignes ;
    \item de calculer explicitement l’inverse d’une matrice ;
    \item d’interpréter géométriquement des volumes orientés ;
    \item d’exprimer certaines propriétés de changement de base.
\end{itemize}

\begin{idee}
Le déterminant mesure, au sens algébrique, comment une matrice déforme l’espace.
\end{idee}

\begin{remarque}
Sur \(\R^2\), le déterminant mesure l’aire orientée du parallélogramme engendré par deux vecteurs.  
Sur \(\R^3\), il mesure le volume orienté du parallélépipède engendré par trois vecteurs.
\end{remarque}

\section{Formes multilinéaires alternées}

\subsection{Multilinéarité}

\begin{definition}
Soit \(E\) un \(\K\)-espace vectoriel.  
Une application
\[
f:E^n\to\K
\]
est dite \textbf{\(n\)-linéaire} si elle est linéaire par rapport à chacune de ses variables lorsque les autres sont fixées.
\end{definition}

\begin{remarque}
Cela signifie que, pour chaque position, si l’on remplace un vecteur par une combinaison linéaire, on peut sortir les coefficients comme dans une application linéaire ordinaire.
\end{remarque}

\begin{exemple}
Une application bilinéaire
\[
f:E\times E\to\K
\]
vérifie
\[
f(\lambda x+\mu y,z)=\lambda f(x,z)+\mu f(y,z)
\]
et
\[
f(x,\lambda y+\mu z)=\lambda f(x,y)+\mu f(x,z).
\]
\end{exemple}

\subsection{Alternance}

\begin{definition}
Une application \(n\)-linéaire
\[
f:E^n\to\K
\]
est dite \textbf{alternée} si elle s’annule dès que deux de ses arguments sont égaux.
\end{definition}

\begin{proposition}
Si \(f\) est alternée, alors échanger deux arguments change le signe :
\[
f(\dots,x_i,\dots,x_j,\dots)=-f(\dots,x_j,\dots,x_i,\dots).
\]
\end{proposition}

\begin{proof}
Considérons l’expression
\[
f(\dots,x_i+x_j,\dots,x_i+x_j,\dots).
\]
Comme \(f\) est alternée, elle vaut \(0\). En développant par multilinéarité, on obtient
\[
f(\dots,x_i,\dots,x_i,\dots)
+
f(\dots,x_i,\dots,x_j,\dots)
+
f(\dots,x_j,\dots,x_i,\dots)
+
f(\dots,x_j,\dots,x_j,\dots)
=0.
\]
Les deux termes extrêmes sont nuls par alternance, donc
\[
f(\dots,x_i,\dots,x_j,\dots)+f(\dots,x_j,\dots,x_i,\dots)=0.
\]
\end{proof}

\begin{corollaire}
Si \(f\) est alternée et si une famille \((x_1,\dots,x_n)\) est liée, alors
\[
f(x_1,\dots,x_n)=0.
\]
\end{corollaire}

\begin{proof}
Si la famille est liée, l’un des vecteurs est combinaison linéaire des autres. En développant par multilinéarité, on obtient une somme de termes dans lesquels deux arguments sont égaux, donc nuls par alternance.
\end{proof}

\section{Définition du déterminant}

\subsection{Point de vue vectoriel}

\begin{theoreme}
Soit \(E\) un espace vectoriel de dimension \(n\), muni d’une base \(\mathcal B=(e_1,\dots,e_n)\).  
Il existe une unique application \(n\)-linéaire alternée
\[
\det_{\mathcal B}:E^n\to \K
\]
telle que
\[
\det_{\mathcal B}(e_1,\dots,e_n)=1.
\]
\end{theoreme}

\begin{proof}
L’existence s’obtient en développant tout \(n\)-uplet dans la base \(\mathcal B\), puis en imposant la multilinéarité et l’alternance.

L’unicité vient du fait que toute application \(n\)-linéaire alternée est entièrement déterminée par sa valeur sur une base, grâce au développement multilineaire.
\end{proof}

\begin{definition}
L’application \(\det_{\mathcal B}\) est appelée le \textbf{déterminant dans la base \(\mathcal B\)}.
\end{definition}

\begin{remarque}
Dans \(E=\K^n\), on travaille en général avec la base canonique, et on note simplement \(\det\).
\end{remarque}

\subsection{Point de vue matriciel}

\begin{definition}
Pour une matrice carrée
\[
A\in M_n(\K),
\]
dont les colonnes sont \(C_1,\dots,C_n\in\K^n\), on définit
\[
\det(A)=\det_{\mathcal E}(C_1,\dots,C_n),
\]
où \(\mathcal E\) est la base canonique de \(\K^n\).
\end{definition}

\begin{remarque}
Le déterminant est donc une fonction
\[
\det:M_n(\K)\to\K.
\]
\end{remarque}

\section{Caractérisation axiomatique}

\begin{theoreme}
Il existe une unique application
\[
\det:M_n(\K)\to\K
\]
vérifiant les propriétés suivantes :
\begin{enumerate}[label=\arabic*)]
    \item elle est multilinéaire par rapport aux colonnes ;
    \item elle est alternée par rapport aux colonnes ;
    \item
    \[
    \det(I_n)=1.
    \]
\end{enumerate}
\end{theoreme}

\begin{proof}
C’est la traduction matricielle du théorème précédent appliqué à \(\K^n\) muni de sa base canonique.
\end{proof}

\begin{remarque}
Cette caractérisation est fondamentale : toutes les propriétés du déterminant en découlent.
\end{remarque}

\section{Premières conséquences}

\subsection{Colonnes égales ou dépendantes}

\begin{proposition}
Si deux colonnes d’une matrice sont égales, alors son déterminant est nul.
\end{proposition}

\begin{proof}
C’est une conséquence directe de l’alternance.
\end{proof}

\begin{proposition}
Si les colonnes d’une matrice sont liées, alors son déterminant est nul.
\end{proposition}

\begin{proof}
Le déterminant est multilinéaire alterné, donc s’annule sur toute famille liée.
\end{proof}

\subsection{Effet d’un échange de colonnes}

\begin{proposition}
Si l’on échange deux colonnes d’une matrice, son déterminant change de signe.
\end{proposition}

\begin{proof}
C’est la propriété générale des applications multilinéaires alternées.
\end{proof}

\subsection{Colonne multipliée par un scalaire}

\begin{proposition}
Si l’on multiplie une colonne d’une matrice par \(\lambda\in\K\), alors le déterminant est multiplié par \(\lambda\).
\end{proposition}

\begin{proof}
Cela découle de la linéarité par rapport à chaque colonne.
\end{proof}

\subsection{Ajout d’un multiple d’une colonne à une autre}

\begin{proposition}
Si l’on ajoute à une colonne un multiple d’une autre colonne, le déterminant ne change pas.
\end{proposition}

\begin{proof}
Par multilinéarité, remplacer \(C_j\) par \(C_j+\lambda C_i\) remplace le déterminant par
\[
\det(\dots,C_j,\dots)+\lambda \det(\dots,C_i,\dots).
\]
Le second terme est nul car deux colonnes y sont égales.
\end{proof}

\subsection{Version lignes}

\begin{remarque}
Toutes les propriétés précédentes valent également par rapport aux lignes, car
\[
\det(A)=\det(A^T).
\]
Nous le démontrerons plus loin.
\end{remarque}

\section{Déterminant des matrices triangulaires}

\begin{theoreme}
Si \(A=(a_{ij})\in M_n(\K)\) est triangulaire supérieure ou inférieure, alors
\[
\det(A)=a_{11}a_{22}\cdots a_{nn}.
\]
\end{theoreme}

\begin{proof}
Traitons le cas triangulaire supérieure.

Par multilinéarité, on peut écrire chaque colonne comme une combinaison linéaire des vecteurs de base. Lorsqu’on développe, tout terme contenant deux vecteurs de base égaux s’annule par alternance. Le seul terme non nul possible est celui qui sélectionne sur chaque colonne le coefficient diagonal, ce qui donne
\[
a_{11}a_{22}\cdots a_{nn}\det(I_n)=a_{11}\cdots a_{nn}.
\]
\end{proof}

\begin{corollaire}
Pour une matrice diagonale
\[
A=\mathrm{diag}(a_1,\dots,a_n),
\]
on a
\[
\det(A)=a_1\cdots a_n.
\]
\end{corollaire}

\begin{exemple}
\[
\det\begin{pmatrix}
2 & 1 & 3\\
0 & -1 & 4\\
0 & 0 & 5
\end{pmatrix}
=2\times(-1)\times 5=-10.
\]
\end{exemple}

\section{Déterminant en petite taille}

\subsection{Ordre \(1\)}

\[
\det(a)=a.
\]

\subsection{Ordre \(2\)}

\begin{proposition}
Pour
\[
A=
\begin{pmatrix}
a & b\\
c & d
\end{pmatrix},
\]
on a
\[
\det(A)=ad-bc.
\]
\end{proposition}

\begin{proof}
Par multilinéarité et alternance, c’est la seule formule possible normalisée par \(\det(I_2)=1\).
\end{proof}

\subsection{Ordre \(3\)}

\begin{proposition}
Pour
\[
A=
\begin{pmatrix}
a & b & c\\
d & e & f\\
g & h & i
\end{pmatrix},
\]
on a
\[
\det(A)=aei+bfg+cdh-ceg-bdi-afh.
\]
\end{proposition}

\begin{remarque}
On retrouve ici la règle de Sarrus, qui n’est valable qu’en dimension \(3\).
\end{remarque}

\section{Formule de Leibniz}

\subsection{Permutations}

\begin{definition}
Une \textbf{permutation} de \(\{1,\dots,n\}\) est une bijection
\[
\sigma:\{1,\dots,n\}\to\{1,\dots,n\}.
\]
On note \(\mathfrak S_n\) l’ensemble des permutations de \(\{1,\dots,n\}\).
\end{definition}

\begin{definition}
À toute permutation \(\sigma\in\mathfrak S_n\), on associe sa \textbf{signature}, notée \(\varepsilon(\sigma)\in\{-1,1\}\), qui vaut \(1\) si \(\sigma\) est paire, et \(-1\) si elle est impaire.
\end{definition}

\subsection{Formule}

\begin{theoreme}[Formule de Leibniz]
Pour toute matrice \(A=(a_{ij})\in M_n(\K)\),
\[
\det(A)=\sum_{\sigma\in\mathfrak S_n}\varepsilon(\sigma)\,a_{1\sigma(1)}a_{2\sigma(2)}\cdots a_{n\sigma(n)}.
\]
\end{theoreme}

\begin{proof}
La formule de Leibniz définit une application multilinéaire alternée des colonnes de \(A\), normalisée par \(\det(I_n)=1\). Par unicité du déterminant, elle coïncide avec lui.
\end{proof}

\begin{remarque}
Cette formule est théoriquement fondamentale, mais peu pratique pour des calculs explicites lorsque \(n\) est grand.
\end{remarque}

\section{Développement selon une ligne ou une colonne}

\subsection{Mineurs et cofacteurs}

\begin{definition}
Soit \(A=(a_{ij})\in M_n(\K)\).  
Pour \(1\le i,j\le n\), on note \(A_{ij}\) la matrice obtenue en supprimant la \(i\)-ème ligne et la \(j\)-ème colonne de \(A\).

Le scalaire
\[
M_{ij}=\det(A_{ij})
\]
est appelé le \textbf{mineur} d’indice \((i,j)\).

Le \textbf{cofacteur} d’indice \((i,j)\) est
\[
C_{ij}=(-1)^{i+j}M_{ij}.
\]
\end{definition}

\subsection{Développement de Laplace}

\begin{theoreme}
Pour toute matrice \(A=(a_{ij})\in M_n(\K)\), on a, pour toute ligne \(i\),
\[
\det(A)=\sum_{j=1}^n a_{ij}C_{ij}.
\]
De même, pour toute colonne \(j\),
\[
\det(A)=\sum_{i=1}^n a_{ij}C_{ij}.
\]
\end{theoreme}

\begin{proof}
On démontre la formule en développant le déterminant par multilinéarité selon la ligne ou la colonne choisie. Les signes viennent de l’effet des permutations nécessaires pour ramener la ligne ou la colonne en tête.
\end{proof}

\begin{exemple}
Pour
\[
A=
\begin{pmatrix}
1 & 2 & 0\\
3 & -1 & 4\\
0 & 5 & 2
\end{pmatrix},
\]
développons selon la première ligne :
\[
\det(A)=1
\begin{vmatrix}
-1 & 4\\
5 & 2
\end{vmatrix}
-2
\begin{vmatrix}
3 & 4\\
0 & 2
\end{vmatrix}
+0.
\]
Donc
\[
\det(A)=1((-1)\cdot 2-4\cdot 5)-2(3\cdot 2-0)
=-22-12=-34.
\]
\end{exemple}

\section{Effet des opérations élémentaires}

\subsection{Sur les lignes}

\begin{proposition}
Pour une matrice \(A\in M_n(\K)\) :
\begin{enumerate}[label=\arabic*)]
    \item échanger deux lignes multiplie le déterminant par \(-1\) ;
    \item multiplier une ligne par \(\lambda\) multiplie le déterminant par \(\lambda\) ;
    \item ajouter à une ligne un multiple d’une autre ne change pas le déterminant.
\end{enumerate}
\end{proposition}

\begin{proof}
Ces résultats sont les versions lignes des propriétés déjà établies sur les colonnes, via la transposition ou directement par la définition multilineaire en lignes.
\end{proof}

\subsection{Conséquence pratique}

\begin{methode}
Pour calculer un déterminant :
\begin{enumerate}
    \item on utilise des opérations élémentaires pour simplifier la matrice ;
    \item on suit l’effet de ces opérations sur le déterminant ;
    \item on se ramène idéalement à une matrice triangulaire.
\end{enumerate}
\end{methode}

\begin{exemple}
Calculons
\[
\det\begin{pmatrix}
1 & 2 & 3\\
2 & 4 & 5\\
1 & 1 & 0
\end{pmatrix}.
\]
On effectue
\[
L_2\leftarrow L_2-2L_1,\qquad L_3\leftarrow L_3-L_1.
\]
On obtient
\[
\begin{pmatrix}
1 & 2 & 3\\
0 & 0 & -1\\
0 & -1 & -3
\end{pmatrix}.
\]
Puis on échange \(L_2\) et \(L_3\), ce qui change le signe :
\[
-\det\begin{pmatrix}
1 & 2 & 3\\
0 & -1 & -3\\
0 & 0 & -1
\end{pmatrix}.
\]
Le déterminant vaut donc
\[
-(1\times (-1)\times (-1))=-1.
\]
\end{exemple}

\section{Déterminant d’un produit}

\begin{theoreme}
Pour toutes matrices \(A,B\in M_n(\K)\),
\[
\det(AB)=\det(A)\det(B).
\]
\end{theoreme}

\begin{proof}
Fixons \(B\). L’application
\[
A\mapsto \det(AB)
\]
est multilinéaire alternée par rapport aux colonnes de \(A\). Elle vaut \(\det(B)\) sur \(I_n\). Par unicité du déterminant,
\[
\det(AB)=\det(A)\det(B).
\]
\end{proof}

\begin{corollaire}
Pour tout \(A\in M_n(\K)\) et tout entier \(k\in\N\),
\[
\det(A^k)=\det(A)^k.
\]
\end{corollaire}

\begin{proof}
Par récurrence ou par application répétée du théorème précédent.
\end{proof}

\section{Transposition et déterminant}

\begin{proposition}
Pour toute matrice \(A\in M_n(\K)\),
\[
\det(A^T)=\det(A).
\]
\end{proposition}

\begin{proof}
La formule de Leibniz est symétrique en lignes et en colonnes. On peut aussi remarquer que le déterminant comme fonction des lignes vérifie les mêmes axiomes que comme fonction des colonnes.
\end{proof}

\section{Déterminant et inversibilité}

\begin{theoreme}
Pour \(A\in M_n(\K)\), les propriétés suivantes sont équivalentes :
\begin{itemize}
    \item \(A\) est inversible ;
    \item \(\det(A)\neq 0\).
\end{itemize}
\end{theoreme}

\begin{proof}
Si \(A\) est inversible, alors
\[
I_n=AA^{-1},
\]
donc
\[
1=\det(I_n)=\det(A)\det(A^{-1}).
\]
Ainsi \(\det(A)\neq 0\).

Réciproquement, si \(\det(A)\neq 0\), on montrera plus loin que
\[
A^{-1}=\frac{1}{\det(A)}\,{}^t\!\mathrm{Com}(A),
\]
ce qui prouve que \(A\) est inversible.
\end{proof}

\begin{corollaire}
Une matrice est non inversible si et seulement si son déterminant est nul.
\end{corollaire}

\begin{corollaire}
Les colonnes d’une matrice \(A\in M_n(\K)\) forment une base de \(\K^n\) si et seulement si
\[
\det(A)\neq 0.
\]
\end{corollaire}

\begin{proof}
Les colonnes forment une base si et seulement si \(A\) est inversible.
\end{proof}

\section{Comatrice et formule de l’inverse}

\subsection{Comatrice}

\begin{definition}
La \textbf{comatrice} de \(A\in M_n(\K)\), notée \(\mathrm{Com}(A)\), est la matrice
\[
\mathrm{Com}(A)=(C_{ij})_{1\le i,j\le n},
\]
où \(C_{ij}\) est le cofacteur d’indice \((i,j)\).
\end{definition}

\begin{remarque}
La matrice utile dans la formule de l’inverse sera la transposée de la comatrice.
\end{remarque}

\subsection{Formule fondamentale}

\begin{theoreme}
Pour toute matrice \(A\in M_n(\K)\),
\[
A\,{}^t\!\mathrm{Com}(A)=\det(A)I_n
\qquad \text{et} \qquad
{}^t\!\mathrm{Com}(A)\,A=\det(A)I_n.
\]
\end{theoreme}

\begin{proof}
Regardons le coefficient \((i,j)\) de
\[
A\,{}^t\!\mathrm{Com}(A).
\]
Il vaut
\[
\sum_{k=1}^n a_{ik}C_{jk}.
\]
Si \(i=j\), c’est exactement le développement de \(\det(A)\) selon la ligne \(i\), donc on obtient \(\det(A)\).

Si \(i\neq j\), cette somme est le déterminant d’une matrice ayant deux lignes égales, donc vaut \(0\).

On obtient donc
\[
A\,{}^t\!\mathrm{Com}(A)=\det(A)I_n.
\]
La seconde formule se prouve de la même manière.
\end{proof}

\subsection{Formule de l’inverse}

\begin{corollaire}
Si \(\det(A)\neq 0\), alors
\[
A^{-1}=\frac{1}{\det(A)}\,{}^t\!\mathrm{Com}(A).
\]
\end{corollaire}

\begin{proof}
Il suffit de diviser l’égalité précédente par \(\det(A)\).
\end{proof}

\section{Déterminant et rang}

\subsection{Rang maximal}

\begin{proposition}
Pour \(A\in M_n(\K)\),
\[
\det(A)\neq 0 \Longleftrightarrow \rg(A)=n.
\]
\end{proposition}

\begin{proof}
\[
\det(A)\neq 0
\Longleftrightarrow
A \text{ inversible}
\Longleftrightarrow
\rg(A)=n.
\]
\end{proof}

\subsection{Colonnes liées}

\begin{proposition}
Si \(\det(A)=0\), alors les colonnes de \(A\) sont liées.
\end{proposition}

\begin{proof}
Si les colonnes étaient libres, elles formeraient une base de \(\K^n\), donc \(A\) serait inversible, donc \(\det(A)\neq 0\), contradiction.
\end{proof}

\subsection{Réciproque}

\begin{proposition}
Les colonnes de \(A\) sont liées si et seulement si
\[
\det(A)=0.
\]
\end{proposition}

\begin{proof}
Une implication est déjà connue par alternance. L’autre résulte de la proposition précédente.
\end{proof}

\section{Déterminant d’une matrice de changement de base}

\begin{proposition}
Une matrice de passage entre deux bases d’un espace vectoriel de dimension \(n\) a un déterminant non nul.
\end{proposition}

\begin{proof}
Une matrice de passage est inversible, donc son déterminant est non nul.
\end{proof}

\begin{remarque}
Le déterminant d’une matrice de passage mesure l’effet du changement de base sur les volumes orientés.
\end{remarque}

\section{Interprétation géométrique}

\subsection{En dimension \(2\)}

\begin{proposition}
Pour deux vecteurs
\[
u=\binom{x_1}{y_1},
\qquad
v=\binom{x_2}{y_2}
\]
de \(\R^2\),
\[
\det(u,v)=x_1y_2-x_2y_1.
\]
Sa valeur absolue est l’aire du parallélogramme engendré par \(u\) et \(v\).
\end{proposition}

\subsection{En dimension \(3\)}

\begin{proposition}
Pour trois vecteurs \(u,v,w\in\R^3\), la valeur absolue de
\[
\det(u,v,w)
\]
est le volume du parallélépipède engendré par ces trois vecteurs.
\end{proposition}

\begin{remarque}
Le signe du déterminant encode une orientation.
\end{remarque}

\section{Déterminants par blocs}

\begin{proposition}
Si
\[
A=
\begin{pmatrix}
B & *\\
0 & C
\end{pmatrix}
\]
est une matrice triangulaire par blocs, où \(B\) et \(C\) sont carrées, alors
\[
\det(A)=\det(B)\det(C).
\]
\end{proposition}

\begin{proof}
La matrice est triangulaire si on regarde les blocs diagonaux ; son déterminant vaut donc le produit des déterminants de ces blocs.
\end{proof}

\section{Méthodes pratiques de calcul}

\begin{methode}
Pour calculer un déterminant :
\begin{enumerate}
    \item regarder si la matrice est triangulaire ou presque ;
    \item utiliser des opérations élémentaires qui préservent ou contrôlent le déterminant ;
    \item développer selon une ligne ou une colonne contenant beaucoup de zéros ;
    \item factoriser si possible ;
    \item en taille \(2\) ou \(3\), utiliser les formules explicites.
\end{enumerate}
\end{methode}

\begin{methode}
Pour montrer qu’un déterminant est nul, on peut :
\begin{itemize}
    \item montrer que deux lignes ou deux colonnes sont égales ;
    \item montrer qu’une ligne ou une colonne est combinaison linéaire des autres ;
    \item montrer que la matrice n’est pas inversible ;
    \item faire apparaître une dépendance linéaire évidente.
\end{itemize}
\end{methode}

\section{Erreurs classiques}

\begin{itemize}
    \item Confondre déterminant et trace.
    \item Penser que le déterminant d’une somme vérifie
    \[
    \det(A+B)=\det(A)+\det(B),
    \]
    ce qui est faux en général.
    \item Oublier qu’échanger deux lignes change le signe.
    \item Oublier qu’ajouter un multiple d’une ligne à une autre ne change pas le déterminant.
    \item Se tromper entre mineur et cofacteur.
    \item Utiliser la règle de Sarrus en dimension autre que \(3\).
\end{itemize}

\section{Résumé du chapitre}

\begin{itemize}
    \item Le déterminant est l’unique forme \(n\)-linéaire alternée normalisée par
    \[
    \det(I_n)=1.
    \]
    \item Si deux lignes ou deux colonnes sont égales, le déterminant est nul.
    \item Échanger deux lignes ou deux colonnes change le signe du déterminant.
    \item Une matrice triangulaire a pour déterminant le produit des coefficients diagonaux.
    \item On a
    \[
    \det(AB)=\det(A)\det(B),
    \qquad
    \det(A^T)=\det(A).
    \]
    \item Une matrice est inversible si et seulement si son déterminant est non nul.
    \item Si \(\det(A)\neq 0\), alors
    \[
    A^{-1}=\frac{1}{\det(A)}\,{}^t\!\mathrm{Com}(A).
    \]
    \item En dimension \(2\) et \(3\), le déterminant admet une interprétation géométrique en aire/volume orienté.
\end{itemize}

\section*{Exercices d’application immédiate}

\begin{enumerate}[label=\textbf{Exercice \arabic*.}]
    \item Calculer :
    \[
    \det\begin{pmatrix}
    2 & 3\\
    1 & 4
    \end{pmatrix}.
    \]

    \item Calculer :
    \[
    \det\begin{pmatrix}
    1 & 2 & 0\\
    3 & -1 & 4\\
    0 & 5 & 2
    \end{pmatrix}.
    \]

    \item Montrer que si deux colonnes d’une matrice sont égales, alors son déterminant est nul.

    \item Montrer qu’une matrice triangulaire a pour déterminant le produit de ses coefficients diagonaux.

    \item Calculer le déterminant de
    \[
    \begin{pmatrix}
    1 & 2 & 3\\
    2 & 4 & 5\\
    1 & 1 & 0
    \end{pmatrix}
    \]
    par opérations élémentaires.

    \item Montrer que
    \[
    \det(AB)=\det(A)\det(B)
    \]
    dans un cas \(2\times 2\).

    \item Déterminer l’inverse de
    \[
    \begin{pmatrix}
    1 & 2\\
    3 & 4
    \end{pmatrix}
    \]
    à l’aide de la formule du déterminant.

    \item Montrer que les colonnes d’une matrice carrée forment une base si et seulement si son déterminant est non nul.

    \item Interpréter géométriquement
    \[
    \det\begin{pmatrix}
    1 & 2\\
    3 & 4
    \end{pmatrix}.
    \]

    \item Montrer qu’une matrice de passage a un déterminant non nul.
\end{enumerate}

\section{Compléments classiques issus des polycopiés de référence}

\subsection{Déterminant de Vandermonde}

\begin{proposition}
Pour des scalaires \(x_0,\dots,x_n\), on définit le déterminant de Vandermonde par
\[
V(x_0,\dots,x_n)=
\begin{vmatrix}
1&1&\cdots&1\\
x_0&x_1&\cdots&x_n\\
x_0^2&x_1^2&\cdots&x_n^2\\
\vdots&\vdots&&\vdots\\
x_0^n&x_1^n&\cdots&x_n^n
\end{vmatrix}.
\]
Alors
\[
V(x_0,\dots,x_n)=\prod_{0\le j<i\le n}(x_i-x_j).
\]
En particulier, \(V(x_0,\dots,x_n)\neq 0\) si et seulement si les \(x_i\) sont deux à deux distincts.
\end{proposition}

\subsection{Règle de Cramer}

\begin{theoreme}[règle de Cramer]
Soit \(A\in M_n(\K)\) inversible et \(AX=B\).
Si \(\Delta=\det(A)\neq 0\), alors le système admet une unique solution, donnée par
\[
x_k=\frac{\Delta_k}{\Delta},
\]
où \(\Delta_k\) est le déterminant obtenu en remplaçant la \(k\)-ième colonne de \(A\) par le second membre \(B\).
\end{theoreme}

\subsection{Déterminants tridiagonaux}

\begin{definition}
Une matrice \textbf{tridiagonale} est une matrice carrée dont tous les coefficients sont nuls en dehors de la diagonale principale et des deux diagonales voisines.
\end{definition}

\begin{proposition}
Pour une suite de déterminants tridiagonaux \((D_n)\), un développement suivant la dernière ligne ou colonne conduit à une relation de récurrence linéaire d’ordre \(2\) de la forme
\[
D_n=a_nD_{n-1}-b_nD_{n-2}.
\]
Cela relie naturellement calcul de déterminants et suites récurrentes linéaires.
\end{proposition}

\end{document}
